\documentclass[11pt]{article}
\usepackage[margin=1.1in]{geometry}
\usepackage[T1]{fontenc}
\usepackage{lmodern}
\usepackage{amsmath,amssymb,amsthm,mathtools}
\usepackage{microtype}
\usepackage{enumitem}
\usepackage[colorlinks=true,linkcolor=black,citecolor=black,urlcolor=black]{hyperref}

\newtheorem{maintheorem}{Theorem}
\renewcommand{\themaintheorem}{\Alph{maintheorem}}
\newtheorem{lemma}{Lemma}
\newtheorem{proposition}[lemma]{Proposition}
\theoremstyle{definition}
\newtheorem{definition}[lemma]{Definition}
\theoremstyle{remark}
\newtheorem{remark}[lemma]{Remark}

\newcommand{\K}{\mathbb{K}}
\newcommand{\R}{\mathbb{R}}
\newcommand{\C}{\mathbb{C}}
\newcommand{\A}{\mathbb{A}}
\newcommand{\PP}{\mathbb{P}}
\newcommand{\M}{\mathcal{M}}
\newcommand{\cO}{\mathcal{O}}
\newcommand{\fm}{\mathfrak{m}}
\newcommand{\fn}{\mathfrak{n}}
\newcommand{\Mbar}{\overline{\M}}
\newcommand{\din}{d_{\mathrm{in}}}
\newcommand{\dout}{d_{\mathrm{out}}}
\newcommand{\dpar}{d_{\mathrm{par}}}
\newcommand{\sm}{\setminus}
\newcommand{\To}{\Theta^{\circ}}
\DeclareMathOperator{\Span}{Span}
\DeclareMathOperator{\Spec}{Spec}

\title{Subnetworks are the only singularities\\[4pt]\large A converse to Theorem 25 of the Primer}
\author{}
\date{November 2026}

\begin{document}
\maketitle

\noindent
Theorem~25 of the Primer shows that the functions of subnetworks of an MLP with polynomial activation are singular points of the neurovariety $\Mbar$. Here we prove the converse: a function $f_\theta$ computed by a parameter $\theta$ at which no neuron is inactive or redundant is a smooth point of $\Mbar$. For deep networks the statement is conditional on a strengthening of the regularity assumption of Theorem~23 (Definition~\ref{def:strong}); for shallow networks we verify this assumption, and the result is unconditional (Theorem~\ref{thm:shallow}). Whether the assumption holds for deep networks and a generic activation of large degree is left open (Remark~\ref{rem:deep}).

\section{Setting and results}

Throughout, $\K$ is a field of characteristic $0$ with algebraic closure $\overline{\K}$, and $\sigma=\sum_{m=1}^r a_mt^m\in\K[t]$ has degree $r\ge2$ and satisfies $\sigma(0)=0$. As in Theorem~23, we consider an MLP with $L\ge2$ layers, widths $d_0,\dots,d_L$ and vanishing biases,
\[
f_\theta(x)=W_L\,\sigma\big(W_{L-1}\cdots\sigma(W_1x)\cdots\big),\qquad \theta=(W_1,\dots,W_L)\in\Theta=\K^{\dpar},\quad \dpar=\textstyle\sum_{j=1}^Ld_jd_{j-1},
\]
with parametrization map $\varphi\colon\Theta\to V=\K[x_1,\dots,x_{d_0}]_{\le r^{L-1}}^{d_L}$, neuromanifold $\M=\varphi(\Theta)$ and neurovariety $\Mbar$. Let $\Sigma$ be the homogeneity group of $\sigma$ (Example~9). We assume that $\sigma$ is not a monomial, so that $\Sigma$ is a finite group of roots of unity; we write $\sigma(\lambda t)=c_\lambda\sigma(t)$ for $\lambda\in\Sigma$. A neuron is inactive or redundant as in Example~9, and
\[
\To=\{\theta\in\Theta:\ \text{no neuron is inactive or redundant at }\theta\},
\]
a non-empty Zariski-open subset of $\Theta$, since $\Sigma$ is finite. The finite group $G$ generated by the permutations and the rescalings by $\Sigma$ of the hidden neurons (see the Identifiability section of the Primer) acts linearly on $\Theta$, preserves $\To$, and satisfies $f_{g\cdot\theta}=f_\theta$. For $D\ge1$, let
\[
P_D=\{p\in\K[x_1,\dots,x_{d_0}]:\ \deg p\le D,\ p(0)=0\},
\]
which contains the pre-activations of layer $j$ when $D=r^{j-1}$. A family of polynomials is \emph{$\Sigma$-distinct} if its members are nonzero and lie in pairwise distinct $\Sigma$-orbits.

Let $R=\overline{\K}[[s]]$, with fraction field $F=\overline{\K}((s))$. An element of $F\otimes U$, for a $\overline\K$-vector space $U$, is \emph{integral} if it lies in $R\otimes U$; its \emph{reduction} is then its value at $s=0$. Elements of $F\otimes P_D$ are polynomials in $x$ with coefficients in $F$: they are curves $s\mapsto p(s)$ in $P_D$, possibly with a pole at $s=0$.

\begin{definition}[Strong regularity]\label{def:strong}
Let $d,D\ge1$. The activation $\sigma$ is \emph{$d$-strongly regular on $P_D$} if the following hold over $\overline{\K}$, i.e., with $P_D$ denoting the corresponding space of polynomials with coefficients in $\overline\K$.
\begin{enumerate}[label=\textup{(R\arabic*)},leftmargin=3.2em,nosep]
\item\label{R1} \emph{(Points.)} If $p_1,\dots,p_k\in P_D$ are $\Sigma$-distinct and $k\le2d$, then $\sigma\circ p_1,\dots,\sigma\circ p_k$ are linearly independent.
\item\label{R2} \emph{(Tangents.)} If $p_1,\dots,p_d\in P_D$ are $\Sigma$-distinct, $c_i\in\overline\K$, $q_i\in P_D$, and
\[
\textstyle\sum_{i=1}^d\big(c_i\,\sigma\circ p_i+(\sigma'\circ p_i)\,q_i\big)=0,
\]
then $c_i=0$ and $q_i=0$ for all $i$.
\item\label{R3} \emph{(Limits.)} Let $p_1(s),\dots,p_d(s)\in F\otimes P_D$ be $\Sigma$-distinct, and let $c_1(s),\dots,c_d(s)\in F\sm\{0\}$. Assume that $\sum_ic_i(s)\,\sigma\circ p_i(s)$ is integral, with reduction $\sum_{i=1}^dc_i^0\,\sigma\circ p_i^0$, where $c_i^0\in\overline\K\sm\{0\}$ and $p_1^0,\dots,p_d^0\in P_D$ are $\Sigma$-distinct. Then, after a permutation of the indices and after replacing each $p_i(s)$ by $\lambda_ip_i(s)$ for some $\lambda_i\in\Sigma$, every $p_i(s)$ is integral, with reduction $p_i^0$.
\end{enumerate}
\end{definition}

Condition \ref{R1} is the $2d$-regularity of Definition~22, as used in Theorem~23. Conditions \ref{R2} and \ref{R3} are its infinitesimal and asymptotic versions. Let $X\subseteq\K[x]$ be the set of the ``neurons'' $c\,\sigma\circ p$. Then \ref{R1} says that distinct neurons are independent, \ref{R2} that the tangent spaces of $X$ at $d$ distinct neurons are independent and that $p\mapsto\sigma\circ p$ is an immersion, and \ref{R3} that a sum of $d$ distinct neurons is not a limit of sums of $d$ neurons that collide or escape to infinity.

\begin{maintheorem}\label{thm:main}
Assume that $\sigma$ is $d_j$-strongly regular on $P_{r^{j-1}}$ for $j=1,\dots,L-1$. Then, for every $\theta\in\To$, the function $f_\theta$ is a smooth point of $\Mbar$. Moreover $\dim\M=\dpar$ and, if $\K$ is algebraically closed, $\varphi(\To)$ is Zariski open in $\Mbar$. Consequently, every singular point of $\Mbar$ that lies in $\M$ is the function of a subnetwork: it is computed by an MLP with the same activation and one hidden neuron fewer.
\end{maintheorem}

For shallow networks, the assumption can be checked.

\begin{maintheorem}\label{thm:shallow}
Let $L=2$, and assume that $\sigma$ has $2d_1$ consecutive nonzero coefficients: $a_m\ne0$ for $m_0\le m<m_0+2d_1$, for some $m_0\ge1$. For instance, $\sigma$ is generic of degree $r\ge2d_1$ with $\sigma(0)=0$, or $\sigma(t)=(1+t)^r-1$ with $r\ge2d_1$. Then $\Sigma=\{1\}$, $\sigma$ is $d_1$-strongly regular on $P_1$, and the conclusions of Theorem~\ref{thm:main} hold. If moreover $(d_0,d_2)\ne(1,1)$ or $\sigma$ has at least $2d_1+1$ nonzero coefficients, then the characterization is exact: $f\in\M$ is a singular point of $\Mbar$ if and only if $f=f_\theta$ for some $\theta$ with an inactive or redundant neuron, i.e., if and only if $f$ is computed by a network with $d_1-1$ hidden neurons.
\end{maintheorem}

In the excluded case, subnetworks need not be singular: for $d_0=d_1=d_2=1$ and $\sigma(t)=t+t^2$, $\Mbar=\{f\in V:f(0)=0\}\cong\K^2$ is a linear space, with no singular point. For $\sigma(t)=t+t^2+t^3$ and the same widths, $\Mbar=\{(X,Y,Z):XZ=Y^2\}$ in the coordinates $f=Xx+Yx^2+Zx^3$, a quadric cone whose only singular point is $0$, the function of the network with no hidden neuron.

\section{A smoothness criterion}

The proof of Theorem~\ref{thm:main} rests on the following criterion, which isolates the three properties of the parametrization that are needed: its fibers are orbits, it is an immersion, and it is proper near the point.

\begin{lemma}\label{lem:criterion}
Let $\K$ be algebraically closed of characteristic $0$, $R=\K[[s]]$ and $F=\K((s))$. Let $X\subseteq\A^M$ be an irreducible closed subvariety, $\varphi\colon X\to\A^N$ a morphism, and $Z$ the Zariski closure of $\varphi(X)$. Let a finite group $G$ act on $X$ with $\varphi\circ g=\varphi$ for all $g\in G$, and let $X^\circ\subseteq X$ be a non-empty $G$-stable open subset consisting of smooth points of $X$. Assume that, for every $x\in X^\circ$:
\begin{enumerate}[label=\textup{(\alph*)},nosep]
\item $\varphi^{-1}(\varphi(x))=G\cdot x$;
\item the differential $d\varphi_x\colon T_xX\to T_{\varphi(x)}\A^N$ is injective;
\item if $\gamma\in X^\circ(F)$ is such that $\varphi(\gamma)\in R^N$ and $\varphi(\gamma)|_{s=0}=\varphi(x)$, then $\gamma\in R^M$.
\end{enumerate}
Then every point of $\varphi(X^\circ)$ is a smooth point of $Z$, $\dim Z=\dim X$, and $\varphi(X^\circ)$ is open in $Z$.
\end{lemma}

\begin{proof}
Fix $x_0\in X^\circ$ and $z_0=\varphi(x_0)$.

\emph{Step 1: $\varphi$ is proper over a neighborhood of $z_0$.} Let $\Gamma\subseteq\A^M\times Z$ be the graph of $\varphi$, which is closed and isomorphic to $X$, and let $\overline X$ be its closure in $\PP^M\times Z$. The projection $\overline\varphi\colon\overline X\to Z$ is proper, and $\overline X\cap(\A^M\times Z)=\Gamma$. We claim that $\overline\varphi^{-1}(z_0)\subseteq\Gamma$. Let $b\in\overline X$ with $\overline\varphi(b)=z_0$. The graph $\Gamma^\circ$ of $\varphi|_{X^\circ}$ is open and dense in $\overline X$, so there is an irreducible curve in $\overline X$ through $b$ that meets $\Gamma^\circ$. Completing the local ring of its normalization at a point above $b$ gives a morphism $\Spec R\to\overline X$ that maps the closed point to $b$ and the generic point into $\Gamma^\circ$. Its generic point is a $\gamma\in X^\circ(F)$, and its composition with $\overline\varphi$ is an $R$-point of $Z$ with closed point $z_0$: thus $\varphi(\gamma)\in R^N$ and $\varphi(\gamma)|_{s=0}=z_0$. By (c), $\gamma\in R^M$. Hence the $R$-point of $\PP^M$ extending $\gamma$, which is unique, lies in $\A^M$, so that $b\in\overline X\cap(\A^M\times Z)=\Gamma$. This proves the claim. Therefore $z_0$ is not in the closed set $\overline\varphi(\overline X\sm\Gamma)$, and over its complement $U_1\ni z_0$ the map $\varphi^{-1}(U_1)\to U_1$ coincides with $\overline\varphi$, hence is proper and surjective.

\emph{Step 2: $\varphi$ is finite over a neighborhood of $z_0$.} By (a), the fiber over $z_0$ is finite. For a proper morphism, the points of the target with finite fiber form an open set, by the semicontinuity of the dimension of the fibers. Hence there is an affine open $U\subseteq U_1$ containing $z_0$ over which $\varphi$ is proper with finite fibers, hence finite. Let $B=\cO(U)$ and $A=\cO(\varphi^{-1}(U))$. Then $B\subseteq A$ is a finite extension of domains, $G$ acts on $A$ by $B$-algebra automorphisms, and $U\subseteq\varphi(X)$.

\emph{Step 3: $\widehat{\cO}_{Z,z_0}\cong\widehat{\cO}_{X,x_0}$.} Let $\fm\subset B$ be the maximal ideal of $z_0$. By (a), the maximal ideals of $A$ above $\fm$ are the ideals $\fn_y$ of the points $y\in G\cdot x_0\subseteq X^\circ$. Since $A$ is finite over $B$,
\[
A\otimes_B\widehat B_\fm\;\cong\;\prod_{y\in G\cdot x_0}\widehat A_{\fn_y},
\]
and $\widehat B_\fm\to A\otimes_B\widehat B_\fm$ is injective, because $\widehat B_\fm$ is flat over $B$. An element $g\in G$ induces an isomorphism of $\widehat B_\fm$-algebras $\widehat A_{\fn_{x_0}}\to\widehat A_{\fn_{g^{-1}x_0}}$. So an element of $\widehat B_\fm$ that vanishes in $\widehat A_{\fn_{x_0}}$ vanishes in every factor, hence is zero: the map $\widehat B_\fm\to\widehat A_{\fn_{x_0}}$ is injective. It is also surjective. Indeed, by (b), the dual map $\fm/\fm^2\to\fn_{x_0}/\fn_{x_0}^2$ of $d\varphi_{x_0}$ is surjective, so that $\fn_{x_0}=\fm A_{\fn_{x_0}}+\fn_{x_0}^2$ and, by Nakayama's lemma, $\fm$ generates the maximal ideal of $\widehat A_{\fn_{x_0}}$. Since the residue fields are both $\K$, a local homomorphism of complete Noetherian local rings with this property is surjective. Thus $\widehat{\cO}_{Z,z_0}=\widehat B_\fm\cong\widehat A_{\fn_{x_0}}=\widehat{\cO}_{X,x_0}$.

Since $x_0$ is a smooth point of $X$, the ring $\widehat{\cO}_{X,x_0}$ is regular of dimension $\dim X$. Hence so is $\cO_{Z,z_0}$: $z_0$ is a smooth point of $Z$, and $\dim Z=\dim X$. Finally, $\varphi(X^\circ)$ is open in $Z$: with $U$ as in Step~2, the set $\varphi(\varphi^{-1}(U)\sm X^\circ)$ is closed in $U$, since finite morphisms are closed, and does not contain $z_0$ by (a). Its complement in $U$ is an open neighborhood of $z_0$ contained in $\varphi(X^\circ)$, since $U\subseteq\varphi(X)$.
\end{proof}

All three assumptions are needed. The cuspidal cubic $t\mapsto(t^2,t^3)$ satisfies (a) and (c), but not (b) at $t=0$. The nodal cubic $t\mapsto(t^2-1,t(t^2-1))$, restricted to $t\ne-1$, satisfies (a) and (b), but not (c) at $t=1$: a second branch through the node comes from the missing point. In our setting, the missing points are the parameters at infinity and the parameters of subnetworks, and \ref{R3} excludes that $f_\theta$, for $\theta\in\To$, is approached from them.

\section{Proof of Theorem~\ref{thm:main}}

\paragraph{Reduction to $\K=\overline\K$.} A polynomial $g$ on $V$ vanishes on $\M$ if and only if $g\circ\varphi=0$ in $\K[\theta]$, since $\K$ is infinite. So the ideal $I$ of $\Mbar$ is the kernel of the $\K$-linear map $g\mapsto g\circ\varphi$, and $I\otimes_\K\overline\K$ is the kernel of the same map over $\overline\K$, i.e., the ideal of the neurovariety of the model over $\overline\K$. Hence generators of $I$ generate the latter, the two neurovarieties have the same dimension, and the Jacobian criterion at a point of $V$ gives the same answer over $\K$ and over $\overline\K$. Moreover, a neuron is inactive or redundant at $\theta\in\K^{\dpar}$ over $\K$ if and only if it is so over $\overline\K$: if two nonzero rows with entries in $\K$ differ by a factor $\lambda$, then $\lambda\in\K$. We therefore assume that $\K$ is algebraically closed. This also covers $\K=\R$, where the singularities of $\M$ are defined through the complex model.

We apply Lemma~\ref{lem:criterion} to $X=\Theta=\A^{\dpar}$, the map $\varphi$, the group $G$ and $X^\circ=\To$, and verify (a), (b) and (c) by induction on $L$. For a parameter $\theta=(W_1,\dots,W_L)$, let $\theta_{<L}=(W_1,\dots,W_{L-1})$, a parameter of the MLP with $L-1$ layers and widths $d_0,\dots,d_{L-1}$, whose function $p_\theta=(p_1,\dots,p_{d_{L-1}})$ is the vector of pre-activations of layer $L-1$. Its coordinates lie in $P_D$, with $D=r^{L-2}$. Writing $d=d_{L-1}$ and $v_i=W_L[:,i]$,
\begin{equation}\label{eq:last}
f_\theta=\sum_{i=1}^{d}v_i\;\sigma\circ p_i.
\end{equation}
The assumption of the theorem holds for the truncated network. If $\theta\in\To$, then $\theta_{<L}$ has no inactive or redundant neuron either, since the conditions on the neurons of layers $j\le L-2$ only involve $W_1,\dots,W_{L-1}$. The same holds after permuting or rescaling the neurons of layer $L-1$, i.e., the rows of $W_{L-1}$. For $L=1$, the map $\varphi$ is a linear isomorphism, $G$ is trivial, and (a), (b), (c) are clear. Let $L\ge2$.

\begin{lemma}\label{lem:preact}
Let $E$ be $\K$ or $F$. For every $\theta\in\To(E)$, the pre-activations $p_1,\dots,p_d$ of layer $L-1$ are $\Sigma$-distinct, and $v_i\ne0$ for all $i$.
\end{lemma}

\begin{proof}
The $v_i$ are nonzero since no neuron of layer $L-1$ is inactive. Let $E=\K$. By induction on $j<L$, the coordinates of the vector $h^{(j)}$ of activations of layer $j$ are linearly independent, where $h^{(0)}=x$: the pre-activations $W_{j+1}h^{(j)}$ of layer $j+1$ are then nonzero and lie in distinct $\Sigma$-orbits, because the rows of $W_{j+1}$ do, and their images under $\sigma$ are linearly independent by \ref{R1} on $P_{r^j}$. Let now $E=F$. The parameters at which $p_1,\dots,p_d$ are $\Sigma$-distinct form an open subset of $\Theta$, defined by polynomial inequations with coefficients in $\K$. By the first case it contains all the $\K$-points of $\To$. By the Nullstellensatz it contains $\To$, hence its $F$-points.
\end{proof}

\paragraph{(a) Fibers.} This is Theorem~23, whose proof uses \ref{R1} one layer at a time. Let $\theta\in\To$ and $f_{\theta'}=f_\theta$, and let $q_i$, $v'_i$ be the pre-activations and the outgoing weights of layer $L-1$ at $\theta'$. By \eqref{eq:last}, $\sum_iv_i\,\sigma\circ p_i=\sum_iv'_i\,\sigma\circ q_i$, where the $p_i$ are $\Sigma$-distinct and $v_i\ne0$ by Lemma~\ref{lem:preact}. Discard the terms with $q_i=0$ and combine those in the same $\Sigma$-orbit, using $\sigma\circ(\lambda q)=c_\lambda\,\sigma\circ q$. By \ref{R1}, applied to each coordinate of the resulting relation among at most $2d$ polynomials, every $p_i$ lies in the orbit of some $q_{\pi(i)}$, with $\pi$ injective, hence a permutation. So no terms were discarded or combined, and, after acting on $\theta'$ by an element of $G$, we get $q_i=p_i$ and $v'_i=v_i$ for all $i$. Then $\theta'_{<L}$ and $\theta_{<L}$ compute the same function $p_\theta$, and by induction they differ by an element of $G$.

\paragraph{(b) Differential.} Let $\theta\in\To$ and $\dot\theta=(\dot W_1,\dots,\dot W_L)$ with $d\varphi_\theta(\dot\theta)=0$. Differentiating \eqref{eq:last},
\[
0=\sum_{i=1}^d\dot v_i\;\sigma\circ p_i+\sum_{i=1}^dv_i\,(\sigma'\circ p_i)\,\dot p_i,
\]
where $\dot v_i=\dot W_L[:,i]$ and $\dot p=(\dot p_1,\dots,\dot p_d)$ is the image of $\dot\theta_{<L}$ under the differential at $\theta_{<L}$ of the parametrization of the truncated network; in particular $\dot p_i\in P_D$. For each output coordinate $k$, \ref{R2} applied with $c_i=\dot v_i[k]$ and $q_i=v_i[k]\,\dot p_i$ gives $\dot v_i[k]=0$ and $v_i[k]\,\dot p_i=0$. Since $v_i\ne0$, we get $\dot W_L=0$ and $\dot p=0$. By induction, $\dot\theta_{<L}=0$.

\paragraph{(c) Properness.} Let $\theta^0\in\To$, and let $\theta(s)\in\To(F)$ be such that $f_{\theta(s)}$ is integral with reduction $f_{\theta^0}$. We show that $\theta(s)$ is integral. Let $p_i(s)\in F\otimes P_D$ and $v_i(s)\in F^{d_L}$ be the pre-activations and outgoing weights of layer $L-1$ at $\theta(s)$, and $p_i^0$, $v_i^0$ those at $\theta^0$. By Lemma~\ref{lem:preact}, both families of pre-activations are $\Sigma$-distinct, and all the $v_i(s)$ and $v_i^0$ are nonzero. Choose $\lambda\in\K^{d_L}$ such that $c_i(s)=\lambda\cdot v_i(s)$ and $c_i^0=\lambda\cdot v_i^0$ are nonzero for all $i$; this excludes finitely many proper $\K$-linear subspaces of $\K^{d_L}$. By \eqref{eq:last}, $\sum_ic_i(s)\,\sigma\circ p_i(s)=\lambda\cdot f_{\theta(s)}$ is integral, with reduction $\sum_ic_i^0\,\sigma\circ p_i^0$. By \ref{R3}, after acting on $\theta(s)$ by an element of $G$, which changes neither $f_{\theta(s)}$ nor the integrality of $\theta(s)$, every $p_i(s)$ is integral with reduction $p_i^0$. Then the $\sigma\circ p_i(s)$ are integral, and their reductions $\sigma\circ p_i^0$ are linearly independent by \ref{R1}. Hence the $v_i(s)$ are integral: otherwise, multiplying \eqref{eq:last} by the least power $s^k$, $k>0$, that makes all of them integral and reducing would give a nontrivial relation among the $\sigma\circ p_i^0$. Thus $W_L(s)$ is integral. Finally, $\theta_{<L}(s)$ has no inactive or redundant neuron, and its function $p_{\theta(s)}$ is integral with reduction $p_{\theta^0}$, the function of $\theta^0_{<L}$. By induction, $\theta_{<L}(s)$ is integral.

\paragraph{Conclusion.} By Lemma~\ref{lem:criterion}, every $f_\theta$ with $\theta\in\To$ is a smooth point of $\Mbar$, $\dim\Mbar=\dpar$, and $\varphi(\To)$ is open in $\Mbar$. Let $f\in\M$ be a singular point of $\Mbar$, and $f=f_\theta$. Then $\theta\notin\To$, so some hidden neuron is inactive or redundant at $\theta$, and it can be removed without changing $f_\theta$, as in Example~9: an inactive neuron contributes $0$, since $\sigma(0)=0$, and a neuron that is redundant with respect to the $k$-th one, with factor $\lambda\in\Sigma$, is absorbed by adding $c_\lambda$ times its outgoing weights to those of the $k$-th one. \qed

\section{Shallow networks: proof of Theorem~\ref{thm:shallow}}

Let $L=2$, $n=d_0$ and $d=d_1$, so that $P_1$ is the space of linear forms $w\cdot x$, $w\in\K^n$, and
\[
\sigma(w\cdot x)=\sum_{m=1}^ra_m\,(w\cdot x)^m
\]
is the decomposition of a neuron into homogeneous components. Hence a polynomial identity among neurons splits into one identity for each $m$ with $a_m\ne0$. Let $\mathcal I=\{m_0,\dots,m_0+2d-1\}$, so that $a_m\ne0$ for $m\in\mathcal I$. Since two consecutive coefficients are nonzero, $\lambda^{m_0}=\lambda^{m_0+1}$ for $\lambda\in\Sigma$, so that $\Sigma=\{1\}$, and $\Sigma$-distinct means distinct and nonzero. As in the previous section, we may assume that $\K$ is algebraically closed. We use the following elementary facts.

\begin{lemma}\label{lem:vdm}
Let $\tau_1,\dots,\tau_k$ be distinct nonzero elements of a field $E$ of characteristic $0$, and let $c_i,\eta_i\in E$.
\begin{enumerate}[label=\textup{(\roman*)},nosep]
\item If $\sum_ic_i\tau_i^m=0$ for $k$ consecutive integers $m\ge0$, then $c_i=0$ for all $i$.
\item If $\sum_i\big(c_i\tau_i^m+m\,\eta_i\tau_i^{m-1}\big)=0$ for $2k$ consecutive integers $m\ge1$, then $c_i=\eta_i=0$ for all $i$.
\end{enumerate}
\end{lemma}

\begin{proof}
Let the integers be $m_1+j$, with $0\le j<k$ in (i) and $0\le j<2k$ in (ii). For $g=\sum_jg_jT^j\in E[T]$ with the same range of $j$, the combination of the identities with coefficients $g_j$ reads $\sum_ic_i\tau_i^{m_1}g(\tau_i)=0$ in (i), and $\sum_i\big(c_i\,h(\tau_i)+\eta_i\,h'(\tau_i)\big)=0$ with $h=T^{m_1}g$ in (ii). In (i), take for $g$ the Lagrange polynomial with $g(\tau_l)=\delta_{il}$. In (ii), take the Hermite polynomials of degree less than $2k$: first $g(\tau_l)=0$ and $g'(\tau_l)=\delta_{il}$ for all $l$, so that $h(\tau_l)=0$ and $h'(\tau_l)=\tau_i^{m_1}\delta_{il}$, which gives $\eta_i=0$; then $g(\tau_l)=\delta_{il}$, which gives $c_i=0$.
\end{proof}

\begin{lemma}\label{lem:generic}
Let $E\supseteq\K$ be a field, and let $u_1,\dots,u_k\in E^n$ be nonzero. Then the $z\in\K^n$ with $u_i\cdot z\ne0$ for all $i$ form the complement of finitely many proper $\K$-linear subspaces of $\K^n$. In particular, they span $\K^n$.
\end{lemma}

\begin{proof}
The map $\K^n\to E$, $z\mapsto u_i\cdot z$, is $\K$-linear and nonzero, since its values at the standard basis vectors are the entries of $u_i$. Its kernel is a proper subspace, and $\K$ is infinite.
\end{proof}

\paragraph{\ref{R1} and \ref{R2}.} Let $w_1,\dots,w_k\in\K^n$ be distinct and nonzero, with $k\le2d$, and $\sum_ic_i\,\sigma(w_i\cdot x)=0$. Then $\sum_ic_i(w_i\cdot x)^m=0$ for all $m\in\mathcal I$. By Lemma~\ref{lem:generic}, applied to the $w_i$ and their differences, there is $z\in\K^n$ such that the $\tau_i=w_i\cdot z$ are distinct and nonzero. Evaluating at $x=z$, Lemma~\ref{lem:vdm}(i) gives $c_i=0$. This is \ref{R1}. For \ref{R2}, let $k=d$ and $q_i=u_i\cdot x$ with $u_i\in\K^n$. The component of degree $m\in\mathcal I$ of the relation in \ref{R2}, divided by $a_m$, is
\[
\textstyle\sum_i\big(c_i\,(w_i\cdot x)^m+m\,(w_i\cdot x)^{m-1}(u_i\cdot x)\big)=0.
\]
Evaluating at $x=z$, for any $z$ as above, Lemma~\ref{lem:vdm}(ii) gives $c_i=0$ and $u_i\cdot z=0$. Since such $z$ span $\K^n$, $u_i=0$.

\paragraph{\ref{R3}.} Let $p_i(s)=w_i(s)\cdot x$ with $w_i(s)\in F^n$ distinct and nonzero, $c_i(s)\in F\sm\{0\}$, and let $p_i^0=w_i^0\cdot x$ and $c_i^0$ be as in \ref{R3}, for $i=1,\dots,d$. Taking homogeneous components, for every $m\in\mathcal I$,
\begin{equation}\label{eq:moments}
\sum_{i=1}^dc_i(s)\,(w_i(s)\cdot x)^m\ \text{ is integral, with reduction }\ \sum_{i=1}^dc_i^0\,(w_i^0\cdot x)^m.
\end{equation}
Let $z\in\K^n$ be such that the $\tau_i^0=w_i^0\cdot z$ are distinct and nonzero, and let $\tau_i(s)=w_i(s)\cdot z\in F$. Evaluating \eqref{eq:moments} at $x=z$, the moments
\[
\mu_j(s)=\sum_ie_i(s)\,\tau_i(s)^j,\qquad e_i=c_i\tau_i^{m_0},\qquad 0\le j\le2d-1,
\]
lie in $R$, with reductions $\mu_j^0=\sum_ie_i^0(\tau_i^0)^j$, where $e_i^0=c_i^0(\tau_i^0)^{m_0}\ne0$. Let $\chi(T)=\prod_i(T-\tau_i(s))=T^d+\sum_{b<d}\chi_bT^b\in F[T]$. Since $\chi(\tau_i)=0$, for $0\le a<d$ we have $\sum_ie_i\tau_i^a\chi(\tau_i)=0$, i.e.,
\[
\sum_{b=0}^{d-1}\mu_{a+b}(s)\,\chi_b=-\mu_{a+d}(s),\qquad 0\le a<d.
\]
The matrix $H(s)=(\mu_{a+b}(s))_{0\le a,b<d}$ of this linear system has entries in $R$, and its reduction $H^0=(\mu_{a+b}^0)$ factors as $\Lambda^\top\operatorname{diag}(e_i^0)\,\Lambda$, where $\Lambda=((\tau_i^0)^b)_{i,b}$ is a Vandermonde matrix. So $\det H^0=\prod_ie_i^0\prod_{i<l}(\tau_i^0-\tau_l^0)^2\ne0$, and $\det H(s)$ is a unit of $R$. By Cramer's rule, $\chi_b\in R$, and the reductions of the $\chi_b$ solve the reduced system, whose unique solution consists of the coefficients of $\prod_i(T-\tau_i^0)$. Thus the $\tau_i(s)$ are roots of a monic polynomial with coefficients in $R$, whose reduction is $\prod_i(T-\tau_i^0)$. Since $R$ is integrally closed, $\tau_i(s)=w_i(s)\cdot z\in R$. By Lemma~\ref{lem:generic}, applied to the $w_i^0$ and their differences, the admissible $z$ span $\K^n$. Hence every $w_i(s)$ is integral, and for every admissible $z$,
\begin{equation}\label{eq:roots}
\prod_{i=1}^d\big(T-w_i(0)\cdot z\big)=\prod_{i=1}^d\big(T-w_i^0\cdot z\big).
\end{equation}
The admissible $z$ are Zariski dense in $\K^n$, so \eqref{eq:roots} is an identity in $\K[T,z_1,\dots,z_n]$. By unique factorization, the factors agree up to order: after a permutation, $w_i(0)=w_i^0$ for all $i$. This is \ref{R3}, and Theorem~\ref{thm:main} applies.

\paragraph{Subnetworks are singular.} It remains to show that, under the last assumption of Theorem~\ref{thm:shallow}, $f_\theta$ is a singular point of $\Mbar$ whenever some neuron is inactive or redundant at $\theta$. This is the analogue of Theorem~25 without biases, and the proof is the same. As in the conclusion of the proof of Theorem~\ref{thm:main}, $f_\theta=f_{\theta'}$ for some $\theta'$ at which the outgoing weights $v_i$ of some neuron $i$ vanish. Replacing the incoming weights of this neuron by any $u\in\K^n$ gives a parameter in the fiber of $f_\theta$, at which the derivative of $\varphi$ with respect to $W_2[k,i]$ is $\sigma(u\cdot x)\,\varepsilon_k$, where $\varepsilon_k$ is the $k$-th standard basis vector of $\K^{d_2}$. The image of the differential of $\varphi$ at any point of the fiber lies in $T_{f_\theta}\Mbar$. Hence $T_{f_\theta}\Mbar\supseteq E^{d_2}$, where $E$ is the span of the $\sigma(u\cdot x)$, $u\in\K^n$. Rescaling $u$ and using a Vandermonde determinant shows that $E$ contains $(u\cdot x)^m$ whenever $a_m\ne0$, and powers of linear forms span the forms of degree $m$. So
\[
\dim T_{f_\theta}\Mbar\;\ge\;d_2\dim E\;=\;d_2\sum_{m:\,a_m\ne0}\binom{n+m-1}{m}.
\]
Each of the $2d$ terms with $m\in\mathcal I$ is at least $n$, with strict inequality if $n\ge2$ and $m\ge2$. If $n\ge2$, this gives $d_2\dim E>2d\,n\,d_2\ge d(n+d_2)$. If $n=1$ and $d_2\ge2$, it gives $d_2\dim E\ge2d\,d_2>d(1+d_2)$. If $n=d_2=1$, then $\dim E$ is the number of nonzero coefficients of $\sigma$, at least $2d+1$ by assumption. In all cases, $\dim T_{f_\theta}\Mbar>d(n+d_2)=\dpar=\dim\Mbar$, and $f_\theta$ is singular. \qed

\section{Remarks}

\begin{remark}[Deep networks]\label{rem:deep}
For $L\ge3$, Theorem~\ref{thm:main} requires strong regularity on $P_D$ with $D=r^{j-1}>1$, where the neurons $\sigma\circ p$ no longer split into homogeneous components. For a generic $\sigma$ of large degree, \ref{R1} is the composition-independence result of Chang and Hendi (2026), conjectured by Kohn et al.\ (2026); it is the assumption of Theorem~23. We do not know whether \ref{R2} and \ref{R3} hold in this generality, and they do not follow formally from \ref{R1}, which only concerns reduced configurations of neurons: \ref{R2} is the same statement for neurons over the dual numbers $\K[\varepsilon]/(\varepsilon^2)$, and \ref{R3} for neurons over $\K((s))$ whose combination is only known modulo $s$. We expect both to hold for generic $\sigma$ of large degree, which by Theorem~\ref{thm:main} and Theorem~25 would characterize the singularities of deep MLPs in $\M$ as the functions of subnetworks.
\end{remark}

\begin{remark}[The assumptions are natural]
Conditions \ref{R2} and \ref{R3} cannot simply be dropped. If $d\varphi_\theta$ is not injective at some $\theta\in\To$, the parametrization ramifies there, as for the cusp. If \ref{R3} fails, some $f_\theta$ with $\theta\in\To$ is also a limit of functions whose parameters degenerate; when these limits form a second branch of $\Mbar$ through $f_\theta$, as for the node, $f_\theta$ is singular. Conversely, a smooth point of $\Mbar$ is unibranch, so that no such second branch exists at the points covered by Theorem~\ref{thm:main}.
\end{remark}

\begin{remark}[Scope]
The theorems concern the points of $\M$. Points of $\Mbar\sm\M$ can be singular as well, and are not functions of the network at all. Monomial activations, for which $\Sigma$ is infinite and the fibers are positive-dimensional, and networks with biases are not covered: the former require replacing $G$ by a torus, and the latter the analogue of Theorem~23 with biases.
\end{remark}

\end{document}
